\chapter{摘\texorpdfstring{\quad}{}要}

涵道风扇无人机因空气动力效率高，安全性强，兼具垂直起降和水平巡航模态等优势，具有广阔的应用前景。然而由于涵道风扇的封闭式旋翼与飞行器周围流场之间存在显著的强耦合效应以及不同飞行模态下的气动特性也不同等原因，使其更容易受到外部扰动的影响。在该因素的影响下，现有涵道风扇无人机的研究在强扰动的影响下不能达到良好的轨迹跟踪效果，这也制约了上层轨迹规划的研究应用。
因此，本文针对涵道风扇无人机的抗扰控制与飞行轨迹规划问题进行研究，聚焦于姿态控制算法、位置控制算法和轨迹规划算法等三个方面，主要研究内容如下：

(1)针对涵道风扇无人机的姿态抗扰控制问题，提出结合增量非线性动态逆和优先级控制分配的姿态控制算法，该算法不依赖准确的系统模型，仅要求对输入通道建模。首先通过模型简化将力矩分为可控力矩和不可控力矩，可控力矩根据其产生机理直接计算，由模型误差和干扰产生的不可控力矩对姿态的影响通过陀螺仪高频采样在每个控制周期内补偿。然后采用优先级控制分配方法将系统力矩输入划分不同优先级，以求最大范围返回容许控制。

(2)针对涵道风扇无人机的位置抗扰控制问题，提出基于增量非线性动态逆的位置控制算法，该算法结合提出的姿态控制形成外环-内环协同的扰动补偿机制。增量非线性动态逆位置控制算法依赖于加速度计的高频采样，通过加速度计测量值估计出机体受到的外力扰动，并在每个控制周期内对外力扰动进行补偿进而得到每个轴的期望加速度。然后基于三轴加速度计算总推力，并且使用一种几何解算方法计算出期望的机体姿态。

(3)针对涵道风扇无人机的飞行轨迹规划问题，提出一种基于快速扩展随机树的改进路径规划算法C-RRT*，C-RRT*算法不仅通过建立两棵随机树彼此交替生长快速找出无碰撞的初始路径，而且还加入了贪婪策略、渐进优化机制和冗余节点删除策略，在提高搜索效率的同时也能提高搜索到的路径质量。针对涵道风扇无人机的约束条件，使用B样条曲线对C-RRT*算法搜索到的初始路径进行平滑处理，并对速度和加速度进行动力学约束得到满足要求的飞行轨迹。

最后，本文通过仿真或实际飞行试验对提出的抗扰控制算法和飞行轨迹规划算法进行验证，结果表明了提出的算法的有效性。

\keywordsCN{涵道风扇无人机；抗扰控制；增量非线性动态逆；飞行轨迹规划}

\chapter{Abstract}

Ducted fan UAVs, characterized by superior aerodynamic performance, high safety, and dual-mode flight capability including vertical takeoff/landing and horizontal cruising, possess broad prospects for practical applications. However, due to the pronounced strong coupling between the shrouded rotor of the ducted fan and the surrounding aerodynamic flow field, as well as the distinct aerodynamic characteristics across different flight modes, the system becomes more susceptible to external disturbances. Under the influence of this factor, existing research on ducted fan UAVs fails to achieve satisfactory trajectory tracking performance in the presence of strong disturbances, which also limits the development and application of high-level trajectory planning. In response to these challenges, this paper investigates the disturbance rejection control and flight trajectory planning of ducted fan UAV, focusing on three key aspects: attitude control algorithms, position control algorithms, and trajectory planning algorithms. The main research contributions are as follows:

(1) To address the attitude disturbance rejection control problem of ducted fan UAVs, an attitude control algorithm combining Incremental Nonlinear Dynamic Inversion and prioritized control allocation is proposed. This approach does not rely on an accurate system model and only requires modeling of the input channels. First, by simplifying the model, the moments are categorized into controllable and uncontrollable components. The controllable moments are directly computed based on their generation mechanisms, while the effects of uncontrollable moments—arising from model uncertainties and external disturbances—are compensated within each control cycle using high-frequency gyroscope sampling. Subsequently, a priority-based control allocation method is employed to assign different priority levels to the system's moment inputs, aiming to maximize the feasible control authority.

(2) To address the position disturbance rejection control problem of ducted fan UAVs, a position control algorithm based on Incremental Nonlinear Dynamic Inversion is proposed. This algorithm, in conjunction with the proposed attitude control method, establishes an outer-loop and inner-loop cooperative disturbance compensation mechanism. The INDI-based position control algorithm relies on high-frequency sampling of accelerometers to estimate external force disturbances acting on the UAV. These disturbances are compensated for within each control cycle to obtain the desired acceleration for each axis. Subsequently, the total thrust is computed based on the three-axis acceleration, and a geometric solving method is employed to determine the desired body attitude.

(3) To address the flight trajectory planning problem of ducted fan UAVs, an improved path planning algorithm, C-RRT*, based on Rapidly-exploring Random Trees, is proposed. The C-RRT* algorithm not only rapidly finds a collision-free initial path by alternately growing two random trees but also incorporates a greedy strategy, progressive optimization mechanism, and redundant node removal strategy to enhance search efficiency while improving path quality. Considering the constraints of the ducted fan UAV, the initial path generated by the C-RRT* algorithm is smoothed using B-spline curves. Additionally, velocity and acceleration constraints are applied to ensure the generated flight trajectory satisfies the dynamic requirements.

Finally, the proposed disturbance rejection control algorithms and flight trajectory planning algorithm are validated through simulation or real-flight tests. The results demonstrate the effectiveness of the proposed algorithms.

\keywordsEN{Ducted Fan UAV; Disturbance Rejection Control; Incremental Nonlinear Dynamic Inversion; Flight Trajectory Planning}